% block3_crank_nicolson.tex
%
% Phase 3, Block 3: Crank-Nicolson, the theta-scheme, and Rannacher
% time-stepping.
%
% Self-contained writeup. Compiles standalone with pdflatex.

\documentclass[11pt,a4paper]{article}

\usepackage[T1]{fontenc}
\usepackage[utf8]{inputenc}
\usepackage{amsmath, amssymb, amsthm}
\usepackage{geometry}
\usepackage{hyperref}
\usepackage{enumitem}

\geometry{margin=1in}

\newtheorem{proposition}{Proposition}[section]
\newtheorem{theorem}[proposition]{Theorem}
\newtheorem{remark}[proposition]{Remark}

\title{Phase 3, Block 3:\\Crank--Nicolson, the $\theta$-scheme,\\and Rannacher time-stepping}
\author{Quant Finance Portfolio}
\date{}

\begin{document}
\maketitle

\section{Setting and motivation}

Blocks 1 (FTCS) and 2 (BTCS) are first-order in time: their local
truncation error is $O(\Delta\tau) + O(\Delta x^2)$. To balance the
spatial and temporal contributions one needs $\Delta\tau \sim
\Delta x^2$, which forces $M \sim N^2$ and total cost $O(N^3)$ for
both schemes (FTCS by CFL, BTCS by accuracy).

This block implements \emph{Crank--Nicolson} (CN), the unique scheme
in the convex family of FTCS and BTCS that is
\emph{second-order in time}: $O(\Delta\tau^2) + O(\Delta x^2)$. With
$M \sim N$, CN achieves $O(\Delta x^2)$ precision at $O(N^2)$ total
cost --- a factor-$N$ improvement over FTCS/BTCS in the same
precision regime. CN is the workhorse of finite-difference option
pricers in production.

The deliverable is a $\theta$-scheme infrastructure
(\texttt{quantlib.theta\_scheme}) that subsumes FTCS, BTCS and CN as
parameter choices, plus the CN pricer with Rannacher time-stepping
(\texttt{quantlib.cn}) to suppress spurious oscillations near the
strike kink.

% =====================================================================
\section{The $\theta$-family}
% =====================================================================

The three schemes of Phase 3 are particular cases of a single
parametric family:
\begin{equation}
\frac{V_j^{n+1} - V_j^n}{\Delta\tau}
\;=\;
\theta\,\mathcal{L}_h V_j^{n+1} \,+\, (1-\theta)\,\mathcal{L}_h V_j^n,
\qquad \theta \in [0, 1],
\label{eq:theta-scheme}
\end{equation}
where $\mathcal{L}_h$ is the discrete spatial operator
\[
\mathcal{L}_h V_j
\;=\; \frac{\sigma^2}{2}\,\frac{V_{j+1} - 2V_j + V_{j-1}}{\Delta x^2}
   + \mu\,\frac{V_{j+1} - V_{j-1}}{2\Delta x}
   - r\,V_j.
\]
The cases of interest are
\begin{center}
\begin{tabular}{ccccc}
\hline
$\theta$ & Name & Type & Stability & Time order \\
\hline
$0$    & FTCS & Explicit & Conditional ($\alpha \le 1/2$) & $O(\Delta\tau)$ \\
$1/2$  & \textbf{Crank--Nicolson} & Implicit & Unconditional & $O(\Delta\tau^2)$ \\
$1$    & BTCS & Implicit & Unconditional & $O(\Delta\tau)$ \\
\hline
\end{tabular}
\end{center}
The order jump at $\theta = 1/2$ is the central observation of this
block. It is the pay-off of the symmetry between the new and old
levels: when the right-hand side is the average of $\mathcal{L}_h V$
at $\tau_n$ and $\tau_{n+1}$, the discrete equation is centred at the
\emph{midpoint} $\tau_{n+1/2}$, and the midpoint rule is one order
more accurate than the endpoint rules.

% =====================================================================
\section{Crank--Nicolson explicitly}
% =====================================================================

Setting $\theta = 1/2$ in \eqref{eq:theta-scheme}, expanding $\mathcal{L}_h$,
and rearranging $V^{n+1}$ to the left, $V^n$ to the right:
\begin{equation}
\boxed{\;
\beta_-\,V_{j-1}^{n+1} + \beta_0\,V_j^{n+1} + \beta_+\,V_{j+1}^{n+1}
\;=\;
\gamma_-\,V_{j-1}^n + \gamma_0\,V_j^n + \gamma_+\,V_{j+1}^n.
\;}
\label{eq:cn-stencil}
\end{equation}
With $\alpha = \frac{\sigma^2}{2}\,\Delta\tau/\Delta x^2$ and
$\nu = \mu\,\Delta\tau/(2\Delta x)$ as before, the LHS coefficients
are
\[
\beta_- = -\tfrac{1}{2}\alpha + \tfrac{1}{2}\nu, \quad
\beta_0 = 1 + \alpha + \tfrac{1}{2}r\Delta\tau, \quad
\beta_+ = -\tfrac{1}{2}\alpha - \tfrac{1}{2}\nu,
\]
and the RHS coefficients are
\[
\gamma_- = +\tfrac{1}{2}\alpha - \tfrac{1}{2}\nu, \quad
\gamma_0 = 1 - \alpha - \tfrac{1}{2}r\Delta\tau, \quad
\gamma_+ = +\tfrac{1}{2}\alpha + \tfrac{1}{2}\nu.
\]
Symmetry: $\beta_\pm = -\tfrac12(\text{the corresponding BTCS coefficient minus 1 on diagonal})$,
$\gamma_\pm = +\tfrac12(\text{the corresponding FTCS coefficient minus 1 on diagonal})$.
CN is, structurally, ``half BTCS plus half FTCS''.

% =====================================================================
\section{Why CN is second-order in time: the midpoint rule}
\label{sec:second-order}
% =====================================================================

The non-trivial claim is that the local truncation error of
\eqref{eq:cn-stencil} is $O(\Delta\tau^2) + O(\Delta x^2)$, despite
the fact that the time difference $(V^{n+1} - V^n)/\Delta\tau$ looks
like the same first-order forward Euler quotient that FTCS uses. The
gain comes from \emph{which time level the equation is centred on}.

Expand $V$ around the midpoint $\tau_{n+1/2} = \tau_n + \Delta\tau/2$:
\begin{align*}
V(\tau_{n+1}) &= V_* + \tfrac{1}{2}\Delta\tau\, V'_* + \tfrac{1}{8}\Delta\tau^2\, V''_*
+ \tfrac{1}{48}\Delta\tau^3\, V'''_* + O(\Delta\tau^4), \\
V(\tau_n)     &= V_* - \tfrac{1}{2}\Delta\tau\, V'_* + \tfrac{1}{8}\Delta\tau^2\, V''_*
- \tfrac{1}{48}\Delta\tau^3\, V'''_* + O(\Delta\tau^4),
\end{align*}
where $V_* := V(\tau_{n+1/2})$, $V'_* := V_\tau(\tau_{n+1/2})$, etc.
Subtracting and dividing by $\Delta\tau$:
\begin{equation}
\frac{V(\tau_{n+1}) - V(\tau_n)}{\Delta\tau}
\;=\; V_\tau(\tau_{n+1/2})
\;+\; \frac{\Delta\tau^2}{24}\,V_{\tau\tau\tau}(\tau_{n+1/2})
\;+\; O(\Delta\tau^4).
\label{eq:midpoint-time}
\end{equation}
The $\Delta\tau$ term \emph{cancels by symmetry}: the quotient is a
second-order approximation to $V_\tau$ at the midpoint.

Analogously, adding the two expansions of $\mathcal{L}_h V$ (one at
$\tau_n$, one at $\tau_{n+1}$) and dividing by 2:
\begin{equation}
\tfrac{1}{2}\!\left[\mathcal{L}_h V(\tau_{n+1}) + \mathcal{L}_h V(\tau_n)\right]
\;=\;
\mathcal{L}_h V(\tau_{n+1/2}) \,+\, O(\Delta\tau^2) \,+\, O(\Delta x^2),
\label{eq:midpoint-space}
\end{equation}
where the $O(\Delta\tau^2)$ comes from the time interpolation (linear
average of the endpoints) and the $O(\Delta x^2)$ is the spatial
truncation of $\mathcal{L}_h$. Combining \eqref{eq:midpoint-time} and
\eqref{eq:midpoint-space}, the CN equation
$(V^{n+1} - V^n)/\Delta\tau = \frac12(\mathcal{L}_h V^{n+1} + \mathcal{L}_h V^n)$
is the discretisation of the exact PDE $V_\tau = \mathcal{L} V$
\emph{at the midpoint}, with truncation
\[
\tau_h = O(\Delta\tau^2) + O(\Delta x^2).
\]
FTCS and BTCS, by contrast, evaluate the right-hand side at $\tau_n$ or
$\tau_{n+1}$ (not the midpoint); the time error there is $O(\Delta\tau)$.

% =====================================================================
\section{The linear system and matrix structure}
% =====================================================================

Equation \eqref{eq:cn-stencil} written for the interior nodes $j =
1, \dots, N-1$ becomes
\[
A\,\mathbf{V}^{n+1}_{\mathrm{int}}
\;=\; B\,\mathbf{V}^n_{\mathrm{int}}
   + \mathbf{b}_{\mathrm{bnd}}^n,
\]
where $A$ and $B$ are tridiagonal $(N-1)\times(N-1)$ matrices with
$(\beta_-, \beta_0, \beta_+)$ and $(\gamma_-, \gamma_0, \gamma_+)$ on
their diagonals, and the boundary vector mixes both time levels:
\begin{equation}
\mathbf{b}_{\mathrm{bnd}}^n
= \begin{pmatrix}
\gamma_-\,V_0^n - \beta_-\,V_0^{n+1} \\
0 \\
\vdots \\
0 \\
\gamma_+\,V_N^n - \beta_+\,V_N^{n+1}
\end{pmatrix}.
\label{eq:bnd-cn}
\end{equation}
Both lateral coefficients of the stencils contribute to the boundary
correction, in contrast to BTCS where only the LHS coefficients
appeared.

The per-step algorithm has three phases:
\begin{enumerate}[leftmargin=2em]
\item Apply the RHS stencil $B$ to $\mathbf{V}^n_{\mathrm{int}}$ as
a three-point stencil sweep ($O(N)$, no need to assemble $B$
explicitly).
\item Add the boundary correction $\mathbf{b}_{\mathrm{bnd}}^n$
($O(1)$: only two endpoints are non-zero).
\item Solve $A\,\mathbf{V}^{n+1}_{\mathrm{int}} = (\text{result})$ by
Thomas with the pre-factored $A$ ($O(N)$).
\end{enumerate}
Total cost per step: $O(N)$. Total cost: $O(N \cdot M)$, the same
as BTCS.

\paragraph{Pre-factoring $A$.} The matrix $A$ is constant in time
(again thanks to the constant-coefficient form of the transformed
PDE). Factoring $A$ once at setup and applying the factor at every
step is the standard pattern, identical to BTCS.

% =====================================================================
\section{Stability of the $\theta$-scheme}
% =====================================================================

The von Neumann amplification factor of \eqref{eq:theta-scheme},
ignoring the reactive term, is
\begin{equation}
g(\xi) = \frac{1 - 4\alpha(1-\theta)\sin^2(\xi/2) + i\,(1-\theta)\beta_C\sin\xi}
              {1 + 4\alpha\theta\sin^2(\xi/2) - i\,\theta\beta_C\sin\xi},
\label{eq:theta-amplification}
\end{equation}
with $\beta_C := \mu\Delta\tau/\Delta x$, $\xi \in [-\pi,\pi]$ the
Fourier wavenumber.

\begin{proposition}[Stability of the $\theta$-scheme]
The $\theta$-scheme is unconditionally von Neumann stable if and only
if $\theta \ge 1/2$. For $\theta < 1/2$, $|g| \le 1$ requires the
conditional bound
\[
\alpha \;\le\; \frac{1}{2(1 - 2\theta)}.
\]
In particular, $\theta = 0$ gives $\alpha \le 1/2$ (the FTCS CFL
condition), and $\theta = 1/2$ is the smallest $\theta$ for which the
bound is vacuous.
\end{proposition}

\begin{proof}[Sketch]
Set $s = \sin^2(\xi/2)$, $u = \sin\xi$, so $u^2 = 4s(1-s)$. Then
$|g|^2 = N(s,u)/D(s,u)$ where both numerator and denominator are
quadratic in $\alpha$. The condition $|g|^2 \le 1$ becomes $D - N \ge
0$, which after expansion is
\[
8\alpha s\,\bigl[(2\theta - 1) + 4\alpha\theta(1-\theta)s\bigr]
\,+\,(\text{terms vanishing when }\beta_C = 0)
\,\ge\, 0.
\]
For $\theta \ge 1/2$, both bracketed terms are non-negative for any
$\alpha \ge 0$, $s \in [0,1]$. For $\theta < 1/2$, the first term is
negative and the bound becomes the stated inequality. (Full algebra
in Strikwerda \S6.3, or \cite[\S4.5]{Smith1985}.) \qed
\end{proof}

CN at $\theta = 1/2$ is therefore \emph{exactly on the boundary}
between conditional and unconditional stability. This is the structural
reason why CN, although stable, has weaker damping properties than
BTCS, as we now examine.

% =====================================================================
\section{The kink-induced oscillation problem}
\label{sec:kink}
% =====================================================================

The amplification factor at the Nyquist wavenumber $\xi = \pi$
(highest representable frequency on the grid, $\sin^2(\pi/2) = 1$,
$\sin\pi = 0$) is
\[
g_\theta(\pi) = \frac{1 - 4\alpha(1-\theta)}{1 + 4\alpha\theta}.
\]
For BTCS ($\theta = 1$) and CN ($\theta = 1/2$):
\[
g_{\mathrm{BTCS}}(\pi) = \frac{1}{1 + 4\alpha}
\xrightarrow{\alpha \to \infty} 0,
\qquad
g_{\mathrm{CN}}(\pi) = \frac{1 - 2\alpha}{1 + 2\alpha}
\xrightarrow{\alpha \to \infty} -1.
\]
\textbf{BTCS damps the Nyquist mode exponentially; CN preserves it
with a sign flip.} In the regime of interest for CN ($\Delta\tau \sim
\Delta x$, so $\alpha \sim 1/\Delta x \to \infty$ as $N \to \infty$),
$|g_{\mathrm{CN}}(\pi)|$ is essentially $1$, and the Nyquist mode is
preserved indefinitely while alternating sign in time.

This is harmless for smooth initial data: the Fourier coefficient of
the Nyquist mode in the initial vector $V^0$ is exponentially small,
and ``preserving exponentially small'' yields nothing visible.

\paragraph{The payoff is not smooth.} A vanilla call payoff
$\phi(x) = K(e^x - 1)^+$ has a kink at $x = 0$: $\phi$ is continuous
but $\phi'$ has a jump. The Fourier expansion of a function with such
a kink has coefficients decaying only as $1/k^2$, not exponentially.
The Nyquist coefficient in $V^0$ is therefore $O(1/N^2)$ --- not
exponentially small.

\paragraph{Consequence.} An $O(1/N^2)$ amplitude at the Nyquist mode,
preserved across $M$ time steps with sign flip $g \approx -1$ at each,
produces a visible oscillation pattern in the computed $V^M$:
\begin{itemize}[leftmargin=2em]
\item Bounded ($|g| \le 1$, so the oscillation does not grow).
\item Persistent (each step preserves the magnitude).
\item Localised near the kink (where the kink-induced mode lives).
\item Alternating sign at consecutive grid nodes (the spatial
signature of $\xi = \pi$).
\end{itemize}
The visible artefact is a sawtooth of small amplitude near the
strike, distinct from the explosive sawtooth of FTCS instability
(Phase~3 numerical lesson 1) by being bounded but visible. It
degrades local accuracy near the strike: the price has $O(\Delta x)$
error there instead of $O(\Delta x^2)$, and the Greeks (gamma in
particular) inherit the oscillation.

% =====================================================================
\section{Rannacher time-stepping}
% =====================================================================

The standard fix, due to Rannacher~\cite{Rannacher1984}, is to
replace the first \emph{two} CN time steps by \emph{four}
half-time-step BTCS steps, each of length $\Delta\tau/2$:
\begin{equation}
\underbrace{\mathrm{BTCS}_{\Delta\tau/2}\circ\mathrm{BTCS}_{\Delta\tau/2}\circ\mathrm{BTCS}_{\Delta\tau/2}\circ\mathrm{BTCS}_{\Delta\tau/2}}_{\text{4 half-steps, total time }2\Delta\tau}
\;\circ\;
\underbrace{\mathrm{CN}_{\Delta\tau} \circ \cdots \circ \mathrm{CN}_{\Delta\tau}}_{M-2\text{ full steps}}.
\label{eq:rannacher}
\end{equation}
\paragraph{Why it works.} Each BTCS half-step damps the Nyquist mode
by a factor $g_{\mathrm{BTCS},\Delta\tau/2}(\pi) = 1/(1 + 2\alpha)$
($\alpha$ is the CN Fourier number; the half-step uses $\alpha/2$,
giving $1/(1 + 2\alpha)$). After four half-steps the Nyquist mode is
multiplied by $1/(1+2\alpha)^4$, which for $\alpha = O(N) \gg 1$ is
$O(1/N^4)$ --- well below the $O(\Delta x^2) = O(1/N^2)$ floor of the
discretisation. By the time CN takes over, the kink-induced
high-frequency modes are damped to round-off, and CN proceeds with
effectively smooth data.

\paragraph{Why two full steps.} Giles \&
Carter~\cite{GilesCarter2005} prove that this is the minimum number
that recovers $O(\Delta\tau^2)$ convergence in $L^2$ globally. For
pointwise $L^\infty$ convergence at the kink itself, four full steps
of Rannacher would be needed; $L^2$ is sufficient for the price at
any fixed spot, recovered by interpolation that smooths over the
kink anyway.

\paragraph{Why BTCS, not finer CN.} The artefact is not a function of
$\Delta\tau$ being too large; it is the absence of damping. CN with
arbitrarily small $\Delta\tau$ still has $|g_{\mathrm{CN}}(\pi)| =
1$ in the limit (the bound is sharp). Only a dissipative scheme can
remove the irregular high-frequency content, and BTCS is the natural
choice because the same Thomas factorisation can be reused with the
appropriate coefficients.

\paragraph{Cost.} Four half-steps of BTCS cost as much as two full CN
steps. For $M \ge 10$ the overhead is below 5\%; for $M \ge 100$
below 0.5\%. Negligible compared to the convergence gain.

% =====================================================================
\section{Implementation outline}
% =====================================================================

The block ships two modules:
\begin{itemize}[leftmargin=2em]
\item \texttt{quantlib.theta\_scheme} (Python) /
\texttt{quant::pde::theta\_scheme\_*} (C++): a generic $\theta$-stepper
parametrised by $\theta$. Internally pre-factors the LHS matrix and
runs a configurable number of time steps. Reusable as scaffolding
for Block 4 (where the LHS solve is replaced by PSOR).
\item \texttt{quantlib.cn} (Python) /
\texttt{quant::pde::cn\_*} (C++): the user-facing CN pricer for
European options, with Rannacher smoothing enabled by default
(\texttt{rannacher\_steps=2}). Setting \texttt{rannacher\_steps=0}
gives vanilla CN, exposed to make the oscillation phenomenon
testable.
\end{itemize}

The signature mirrors BTCS:
\begin{verbatim}
cn_european_call(S, K, r, sigma, T, *, N, M, n_sigma=4.0,
                 rannacher_steps=2)
\end{verbatim}

The $\theta$-scheme module exposes a low-level kernel
\texttt{theta\_march(grid, V0, theta, bc\_lower, bc\_upper)} that the
CN pricer calls twice for the Rannacher protocol: once with
$\theta = 1$ for the initial half-stepped BTCS warm-up, then with
$\theta = 1/2$ for the rest.

\paragraph{No replacement of FTCS/BTCS.} The FTCS and BTCS modules of
Blocks 1 and 2 are unchanged. The $\theta$-scheme infrastructure is a
parallel mechanism intended for CN and any future generalisations.
The previous granular commits remain valid; this block adds, does not
modify.

% =====================================================================
\section{Validation strategy}
% =====================================================================

Six tests certify the new pricer:
\begin{enumerate}[leftmargin=2em]
\item Cross-validation against Black--Scholes closed form on a
parameter grid (the standard battery of Blocks 1--2).
\item \textbf{Quadratic convergence in time:} halving $\Delta\tau$
at fixed (small) $\Delta x$ should reduce the error by a factor of
$\approx 4$, not $\approx 2$ as BTCS would. This is \emph{the} test
of CN.
\item \textbf{Balanced quadratic convergence:} doubling both $N$ and
$M$ should reduce the error by $\approx 4$, vs the $\approx 2$
expected for BTCS in the same configuration.
\item \textbf{Vanilla CN (without Rannacher) has degraded convergence:}
running the same problem with \texttt{rannacher\_steps=0} should show
a noticeably worse convergence ratio. This is the empirical evidence
of the kink-induced oscillation.
\item Reproduction of BTCS via $\theta = 1$ in the $\theta$-scheme:
verifies the generic stepper against the Block~2 pricer.
\item Put-call parity, input validation, etc.
\end{enumerate}

The fourth test is the operational signature of Block~3's numerical
lesson, analogous to the CFL-violation diagnostic of Block~1: a
configuration that produces visibly degraded output, captured as a
passing assertion to prevent regression.

% =====================================================================
\section*{References}
% =====================================================================

\begin{thebibliography}{9}

\bibitem{TavellaRandall2000}
D. Tavella and C. Randall.
\emph{Pricing Financial Instruments: The Finite Difference Method.}
Wiley, 2000. \S3.3 (CN), \S3.5 (the $\theta$-scheme), \S6.3 (Rannacher
and other smoothing techniques).

\bibitem{Seydel2017}
R. Seydel.
\emph{Tools for Computational Finance.}
6th ed., Springer, 2017. \S4.2.3 (CN), \S4.5.4 (kink and Rannacher in
practice).

\bibitem{Strikwerda2004}
J.~C. Strikwerda.
\emph{Finite Difference Schemes and Partial Differential Equations.}
2nd ed., SIAM, 2004. \S6.3 (stability analysis of the $\theta$-method).

\bibitem{Smith1985}
G.~D. Smith.
\emph{Numerical Solution of Partial Differential Equations:
Finite Difference Methods.}
3rd ed., OUP, 1985. \S1.13 (the $\theta$-family of schemes).

\bibitem{Rannacher1984}
R. Rannacher.
\emph{Finite element solution of diffusion problems with irregular data.}
Numerische Mathematik 43, 309--327, 1984. The original paper
introducing the smoothing protocol.

\bibitem{GilesCarter2005}
M.~B. Giles and R. Carter.
\emph{Convergence analysis of Crank-Nicolson and Rannacher
time-marching.}
Journal of Computational Finance 9(4), 89--112, 2005. Sharp results
for option pricing: proves that two Rannacher steps suffice for
$L^2$-recovery of $O(\Delta\tau^2)$ convergence.

\end{thebibliography}

\end{document}
